\section{Conclusions and Future Scope}

The project set out to determine whether diffusion or LoRA inpainting of rare
room types improves room-instance segmentation of architectural floor plans, and
whether that answer depends on the backbone. The measured answer is that it does
not improve segmentation, and that the backbone matters far more than the
augmentation.

Goal 1 was achieved as specified. Fourteen controlled arms were built on
leakage-safe generators that read only frozen training identifiers, each accepted
sample passing automated geometry, pixel-quality and exact-structural-invariance
checks with a recorded decision audit and output hash, and each arm training
under the identical model, schedule and 35-step-per-epoch budget in both the
multiclass and class-agnostic views. The matrix included the non-generative
controls that make the diffusion treatment interpretable, and it was that design
choice --- rather than any positive finding --- that produced the project's most
useful result.

Goal 2 was achieved in the sense the goal statement anticipated: a defensible
answer was obtained, and the answer is a null result. Across 84 Mask R-CNN
training runs, no augmentation arm produced a difference from the unaugmented
baseline that the data supports. Every paired sign test on the three-seed grid
returned $p \ge 0.25$, and $0.25$ is the floor of the test at $n=3$. Nine of the
thirteen augmented arms scored below baseline on the class-agnostic track, the
LoRA conversion worst at $-0.057$. The cross-validation over 21 held-out images
reproduced the same ordering. The single result with a consistent sign is
exemplar copy-paste at $+0.045$ class-agnostic AP, positive in all five folds at
$p = 0.0625$ --- and copy-paste is the control arm that contains no generated
pixels, which is the opposite of the hypothesis under test. It is also the
outcome Ghiasi et al.\ \cite{ghiasi2021copypaste} would have predicted.

Goal 3 changed the framing of the project. SAM 3 without any augmentation
reached $0.713$ class-agnostic and $0.115$ multiclass mask AP, exceeding the best
Mask R-CNN arm by $0.089$ and $0.011$ respectively. The class-agnostic gap is
about two and a half times the largest augmentation delta measured anywhere in
the grid. Under augmentation SAM 3 showed a $+0.032$ multiclass gain from
copy-paste, but on a single seed, inside the $0.002$--$0.018$ seed-noise band
measured on the grid, and accompanied by a class-agnostic loss under all three
augmentations --- so the sign of the copy-paste effect is opposite in the two
backbones. That inconsistency is itself the finding: at this data scale the
augmentation effect is not a stable property of the augmentation.

Goal 4 supplied the diagnosis. Handing a frozen DINOv2-base encoder perfect
ground-truth crops --- no detector error whatsoever --- yielded $0.454$ top-1
accuracy over 19 classes and exactly $0.000$ macro recall over the six
primary-rare classes, in every arm and for both heads. Adding the 171 synthetic
rare crops lowered top-1 accuracy in both heads and left rare recall at zero. The
end-to-end two-stage system reached $0.091$ mask AP, below both fine-tuned
detectors. Combined with the class-agnostic results, this locates the failure
precisely: room \emph{localisation} is largely solved on this dataset at
$0.59$--$0.71$ AP, and room \emph{typing} is not, for the tail or for several head
classes either. No augmentation of the training pixels was going to fix a
classification problem of that shape.

The wider conclusion is a scope statement rather than a defeat. The generative
augmentation literature surveyed in Section~3 reports its gains on benchmarks
with 100k training images where a ``rare'' class still has a few hundred
instances \cite{gupta2019lvis,xie2024mosaicfusion,shin2023fillup}. This project
had 35 training plans, 30 rare training instances and 17 rare test instances. In
that regime the measurement itself is the binding constraint: an AP computed over
one or two instances is a step function, and the experiment can exclude a large
rare-class effect but cannot resolve a small one. The qualitative evidence in
Figures~\ref{fig:qastrategya} and \ref{fig:loraval} suggests the mechanism as
well. A pretrained diffusion model asked to draw a German Versammlungsraum
produces perspective furniture, and the LoRA adapter that fixed the domain
mismatch learned mostly that floor plans are white space --- partly because the
mask construction protects detected dark strokes, and the printed room labels and
symbols that identify a room type are exactly those strokes. The safety
machinery that made the augmentation trustworthy also removed most of what
carried the class signal.

\subsection*{Future scope}

Four directions follow from the measurements rather than from speculation, and
they are ordered by expected value per unit of effort.

\emph{Read the drawing's own text.} Every plan in this dataset prints the room
type inside the room --- ``B\"uro'', ``Technik'', ``Abstellraum''. The vision-only
classifier scored zero recall on all six rare classes while that label sat inside
the crop it was given. Pairing the class-agnostic detector, which already works,
with OCR or a vision-language model over the crop is the obvious next experiment,
and WAFFLE \cite{ganon2025waffle} demonstrates the approach at scale on
heterogeneous floor plans. This is likely to move the multiclass number more than
anything attempted in this project.

\emph{Try a loss-side long-tail treatment before any further data synthesis.} The
project spent its effort on the data and none on the objective. A Seesaw loss
\cite{wang2021seesaw} costs one configuration change, needs no generator, and
attacks the classification imbalance the two-stage experiment identified as the
actual failure. It should be run on the SAM 3 backbone, which is where the
headroom is.

\emph{Fix the semantic gap in generation, or steer instead of adapting.} If
generative augmentation is revisited, the two concrete fixes are already
identified in the literature. Gen-n-Val \cite{huang2025gennval} shows that a
semantic validator, not more generation, is the binding constraint --- our
pipeline could certify structural safety but recorded
\texttt{semantic\_correctness = unverified} for every prompt-based sample. DoGE
\cite{wang2024doge} avoids adapter training entirely by steering a frozen
generator with a domain-gap embedding, which suits a 30-instance tail better than
LoRA fine-tuning does. The revised ``semantic v2'' generator implemented late in
this project --- masking compact interior labels while preserving structural
lines, and seeding generation with a real rare exemplar --- is built but was not
run through the full grid, and it is the cheapest of these to complete. TMI
\cite{song2026tmi} additionally predicts that the prompt-driven and
exemplar-driven families are complementary and that neither works alone, which
matches what was observed here and suggests the combination is worth one run.

\emph{Collect more rare-class test data before running anything else.} This is
the least interesting recommendation and the most important one. Seventeen rare
test instances cannot support the claim the project was designed to make, and no
amount of methodological care removes that ceiling. Roughly a hundred rare test
instances --- reachable by annotating a few dozen further plans from the same
archive --- would make the difference between excluding a large effect and
measuring a real one. Until that exists, the defensible use of this dataset is
class-agnostic room extraction, where both backbones already perform well, with
room typing handled by the printed labels rather than by appearance.
